\documentclass[a4paper,11pt,oneside]{memoir}

% Castellano
\usepackage[spanish]{babel}
\selectlanguage{spanish}
\usepackage[utf8]{inputenc}
\usepackage{placeins}

\RequirePackage{booktabs}
\RequirePackage[table]{xcolor}
\RequirePackage{xtab}
\RequirePackage{multirow}
% Links
\usepackage[colorlinks]{hyperref}
\hypersetup{
	allcolors = {red}
}

% Ecuaciones
\usepackage{amsmath}

% Rutas de fichero / paquete
\newcommand{\ruta}[1]{{\sffamily #1}}

% Párrafos
\nonzeroparskip


% Imagenes
\usepackage{graphicx}
\newcommand{\imagen}[2]{
	\begin{figure}[!h]
		\centering
		\includegraphics[width=0.9\textwidth]{#1}
		\caption{#2}\label{fig:#1}
	\end{figure}
	\FloatBarrier
}

\newcommand{\imagenflotante}[2]{
	\begin{figure}%[!h]
		\centering
		\includegraphics[width=0.9\textwidth]{#1}
		\caption{#2}\label{fig:#1}
	\end{figure}
}



% El comando \figura nos permite insertar figuras comodamente, y utilizando
% siempre el mismo formato. Los parametros son:
% 1 -> Porcentaje del ancho de página que ocupará la figura (de 0 a 1)
% 2 --> Fichero de la imagen
% 3 --> Texto a pie de imagen
% 4 --> Etiqueta (label) para referencias
% 5 --> Opciones que queramos pasarle al \includegraphics
% 6 --> Opciones de posicionamiento a pasarle a \begin{figure}
\newcommand{\figuraConPosicion}[6]{%
  \setlength{\anchoFloat}{#1\textwidth}%
  \addtolength{\anchoFloat}{-4\fboxsep}%
  \setlength{\anchoFigura}{\anchoFloat}%
  \begin{figure}[#6]
    \begin{center}%
      \Ovalbox{%
        \begin{minipage}{\anchoFloat}%
          \begin{center}%
            \includegraphics[width=\anchoFigura,#5]{#2}%
            \caption{#3}%
            \label{#4}%
          \end{center}%
        \end{minipage}
      }%
    \end{center}%
  \end{figure}%
}

%
% Comando para incluir imágenes en formato apaisado (sin marco).
\newcommand{\figuraApaisadaSinMarco}[5]{%
  \begin{figure}%
    \begin{center}%
    \includegraphics[angle=90,height=#1\textheight,#5]{#2}%
    \caption{#3}%
    \label{#4}%
    \end{center}%
  \end{figure}%
}
% Para las tablas
\newcommand{\otoprule}{\midrule [\heavyrulewidth]}
%
% Nuevo comando para tablas pequeñas (menos de una página).
\newcommand{\tablaSmall}[5]{%
 \begin{table}
  \begin{center}
   \rowcolors {2}{gray!35}{}
   \begin{tabular}{#2}
    \toprule
    #4
    \otoprule
    #5
    \bottomrule
   \end{tabular}
   \caption{#1}
   \label{tabla:#3}
  \end{center}
 \end{table}
}

%
% Nuevo comando para tablas pequeñas (menos de una página).
\newcommand{\tablaSmallSinColores}[5]{%
 \begin{table}[H]
  \begin{center}
   \begin{tabular}{#2}
    \toprule
    #4
    \otoprule
    #5
    \bottomrule
   \end{tabular}
   \caption{#1}
   \label{tabla:#3}
  \end{center}
 \end{table}
}

\newcommand{\tablaApaisadaSmall}[5]{%
\begin{landscape}
  \begin{table}
   \begin{center}
    \rowcolors {2}{gray!35}{}
    \begin{tabular}{#2}
     \toprule
     #4
     \otoprule
     #5
     \bottomrule
    \end{tabular}
    \caption{#1}
    \label{tabla:#3}
   \end{center}
  \end{table}
\end{landscape}
}

%
% Nuevo comando para tablas grandes con cabecera y filas alternas coloreadas en gris.
\newcommand{\tabla}[6]{%
  \begin{center}
    \tablefirsthead{
      \toprule
      #5
      \otoprule
    }
    \tablehead{
      \multicolumn{#3}{l}{\small\sl continúa desde la página anterior}\\
      \toprule
      #5
      \otoprule
    }
    \tabletail{
      \hline
      \multicolumn{#3}{r}{\small\sl continúa en la página siguiente}\\
    }
    \tablelasttail{
      \hline
    }
    \bottomcaption{#1}
    \rowcolors {2}{gray!35}{}
    \begin{xtabular}{#2}
      #6
      \bottomrule
    \end{xtabular}
    \label{tabla:#4}
  \end{center}
}

%
% Nuevo comando para tablas grandes con cabecera.
\newcommand{\tablaSinColores}[6]{%
  \begin{center}
    \tablefirsthead{
      \toprule
      #5
      \otoprule
    }
    \tablehead{
      \multicolumn{#3}{l}{\small\sl continúa desde la página anterior}\\
      \toprule
      #5
      \otoprule
    }
    \tabletail{
      \hline
      \multicolumn{#3}{r}{\small\sl continúa en la página siguiente}\\
    }
    \tablelasttail{
      \hline
    }
    \bottomcaption{#1}
    \begin{xtabular}{#2}
      #6
      \bottomrule
    \end{xtabular}
    \label{tabla:#4}
  \end{center}
}

%
% Nuevo comando para tablas grandes sin cabecera.
\newcommand{\tablaSinCabecera}[5]{%
  \begin{center}
    \tablefirsthead{
      \toprule
    }
    \tablehead{
      \multicolumn{#3}{l}{\small\sl continúa desde la página anterior}\\
      \hline
    }
    \tabletail{
      \hline
      \multicolumn{#3}{r}{\small\sl continúa en la página siguiente}\\
    }
    \tablelasttail{
      \hline
    }
    \bottomcaption{#1}
  \begin{xtabular}{#2}
    #5
   \bottomrule
  \end{xtabular}
  \label{tabla:#4}
  \end{center}
}



\definecolor{cgoLight}{HTML}{EEEEEE}
\definecolor{cgoExtralight}{HTML}{FFFFFF}

%
% Nuevo comando para tablas grandes sin cabecera.
\newcommand{\tablaSinCabeceraConBandas}[5]{%
  \begin{center}
    \tablefirsthead{
      \toprule
    }
    \tablehead{
      \multicolumn{#3}{l}{\small\sl continúa desde la página anterior}\\
      \hline
    }
    \tabletail{
      \hline
      \multicolumn{#3}{r}{\small\sl continúa en la página siguiente}\\
    }
    \tablelasttail{
      \hline
    }
    \bottomcaption{#1}
    \rowcolors[]{1}{cgoExtralight}{cgoLight}

  \begin{xtabular}{#2}
    #5
   \bottomrule
  \end{xtabular}
  \label{tabla:#4}
  \end{center}
}




\graphicspath{ {./img/} }

% Capítulos
\chapterstyle{bianchi}
\newcommand{\capitulo}[2]{
	\setcounter{chapter}{#1}
	\setcounter{section}{0}
	\chapter*{#2}
	\addcontentsline{toc}{chapter}{#2}
	\markboth{#2}{#2}
}

% Apéndices
\renewcommand{\appendixname}{Apéndice}
\renewcommand*\cftappendixname{\appendixname}

\newcommand{\apendice}[1]{
	%\renewcommand{\thechapter}{A}
	\chapter{#1}
}

\renewcommand*\cftappendixname{\appendixname\ }

% Formato de portada
\makeatletter
\usepackage{xcolor}
\newcommand{\tutor}[1]{\def\@tutor{#1}}
\newcommand{\course}[1]{\def\@course{#1}}
\definecolor{cpardoBox}{HTML}{E6E6FF}
\def\maketitle{
  \null
  \thispagestyle{empty}
  % Cabecera ----------------
\noindent\includegraphics[width=\textwidth]{cabecera}\vspace{1cm}%
  \vfill
  % Título proyecto y escudo informática ----------------
  \colorbox{cpardoBox}{%
    \begin{minipage}{.8\textwidth}
      \vspace{.5cm}\Large
      \begin{center}
      \textbf{TFG del Grado en Ingeniería Informática}\vspace{.6cm}\\
      \textbf{\LARGE\@title{}}
      \end{center}
      \vspace{.2cm}
    \end{minipage}

  }%
  \hfill\begin{minipage}{.20\textwidth}
    \includegraphics[width=\textwidth]{escudoInfor}
  \end{minipage}
  \vfill
  % Datos de alumno, curso y tutores ------------------
  \begin{center}%
  {%
    \noindent\LARGE
    Presentado por \@author{}\\ 
    en Universidad de Burgos --- \@date{}\\
    Tutor: \@tutor{}\\

  }%
  \end{center}%
  \null
  \cleardoublepage
  }
\makeatother


% Datos de portada
\title{PLGRAM \\Documentación Técnica}
\author{Víctor Renuncio Tobarr}
\tutor{Dr. César Ignacio García Osorio y Álvar Arnaiz Gonzalez}

\date{\today}

\begin{document}

\maketitle



\cleardoublepage



%%%%%%%%%%%%%%%%%%%%%%%%%%%%%%%%%%%%%%%%%%%%%%%%%%%%%%%%%%%%%%%%%%%%%%%%%%%%%%%%%%%%%%%%



\frontmatter


\clearpage

% Indices
\tableofcontents

\clearpage

\listoffigures

\clearpage

%\listoftables

%\clearpage

\mainmatter

\appendix

\apendice{Manuales}

\section{Introducción}

En este anexo se detalla el estudio desde el punto de vista temporal y de la viabilidad del proyecto software.

La planificación es una de las tareas más importantes en el desarrollo de un proyecto software y servirá para determinar objetivos, evaluar la viabilidad del proyecto, priorizar actividades...

En la primera parte del anexo se detallará la planificación temporal del proyecto teniendo en cuenta la metodología ágil que se va a utilizar: $Scrum$. En esta fase se determinarán los elementos que forman el \textit{Product Backlog} y la prioridad de cada uno de ellos.
Debido a la metodología empleada, no se utilizará el clásico diagrama de $GANTT$. En lugar de esto, se definirá el \textit{Product Backlog} y para el seguimiento se utilizará una herramienta de gestión especializada en metodologías ágiles, $VersionOne$, que permite el seguimiento diario de las tareas por parte del equipo de desarrollo.

En la segunda parte se justificará el desarrollo del proyecto desde diversos puntos de vista: viabilidad técnica, legal, económica, etc.

\section{Planificación temporal}
Esta sección detalla el trabajo realizado en cada $sprint$ del proyecto. Tras la primera reunión con el tutor, se decidió que cada $sprint$ duraría una semana y se correspondería con una serie de objetivos planteados y unos resultados a conseguir. 

A continuación se describe semana por semana dichos objetivos y los resultados que se fueron consiguiendo. Se incluye adicionalmente un $sprint$ inicial que transcurre entre la asignación del proyecto y el comienzo oficial.

\subsection{Sprint 0: hasta el 3 marzo}

Esta iteración tiene consideraciones especiales, puesto que desde que se asignaron los proyectos hasta que realmente se comenzó a trabajar semana a semana con el tutor, se realizó una búsqueda de información relacionada con la futura aplicación.

Para ello se leyeron los apuntes de la asignatura \textit{Procesadores del Lenguaje}, se buscó información sobre programas similares y se recabó información sobre los programas que iban a ser de utilidad a la hora de desarrollar el proyecto.

Así, se creó el repositorio en $Github$ y se comenzó a realizar la programación temporal utilizando el software online $VersionOne$. 

\subsection{Sprint 1: 3 marzo - 10 marzo}

La primera reunión con el tutor fue el 3 de marzo. En ella, se habló sobre cómo se podría distribuir mejor la aplicación que se iba a crear. Se decidió entonces optar por dos soluciones: un prototipo de línea de comandos y una aplicación similar al prototipo con una interfaz gráfica accesible.

Planteados los objetivos generales del proyecto, se empezó a trabajar en el proyecto actual. 

Se planteó el primer objetivo, que fue empezar a familiarizarse con las clases de $Burgram$. Para facilitar dicha familiarización  con la estructura de clases de $Burgram$ y de $PLQuiz$, el tutor me recomendó instalar un plugin que de forma automática generase diagramas $UML$ con la arquitectura de las aplicaciones. Este plugin sería: \url{http://www.objectaid.com/class-diagram}.

Otro de los objetivos planteados fue crear un prototipo de línea de comandos, concretamente la idea sería tener una pequeña aplicación en modo texto que:

\begin{enumerate}
\item Tome como argumento el nombre de un archivo con la especificación de una gramática en el formato $bison$.
\item Utilice su contenido para, utilizando las clases adecuadas de $Burgram$, instanciar un objeto \textit{Gramática}.
\item Interrogue el objeto \textit{Gramática} para mostrar el \textit{FIRST} y el \textit{FOLLOW} de los símbolos de la gramática.
\end{enumerate}


\subsection{Sprint 2: 10 marzo - 17 marzo}

El 10 de marzo se realizó una reunión en la que se valoraron los resultados del primer \textit{sprint}. Todos los resultados fueron satisfactorios, se consiguió crear un prototipo en el cual, introduciendo una gramática en formato \textit{bison}, se obtenía de ella el \textit{FIRST} y el \textit{FOLLOW} por pantalla.

En esta reunión se plantearon nuevos objetivos: 

\begin{itemize}
\item Generar la tabla de análisis sintáctico predictivo (\textit{TASP}) en el prototipo.
\item Instalar MiK\TeX{}.
\item Instalar \textit{\TeX{}Maker}.
\item Hacer un uso correcto de los Path.
\end{itemize}


\subsection{Sprint 3: 17 marzo - 31 marzo}

El 17 de marzo se realizó la siguiente reunión en la que se valoraron los resultados del segundo \textit{sprint}. Todos los resultados fueron satisfactorios. Se consiguió generar la \textit{TASP} y se preparó el prototipo para que los path no diesen problemas en los diferentes sistemas operativos.
 
Los objetivos que se plantearon en esta reunión fueron:

\begin{itemize}
\item Introducir parámetros en el prototipo.
\item Investigar sobre diferentes motores de plantillas, como por ejemplo $Mustache$, $Trimou$, $JinJava$, $FreeMaker$ o $Shootout$, y decidir cuál era el óptimo para el proyecto.
\item Generar una plantilla para crear ficheros .XML utilizables por \textit{MOODLE}.
\item Generar una plantilla para crear ficheros .TEX. 
\end{itemize}

Se estableció que este \textit{sprint} sería de dos semanas debido a la Semana Santa.



\subsection{Sprint 4: 31 marzo - 7 abril}

El 31 de marzo se realizó la siguiente reunión en la que se valoraron los resultados del tercer \textit{sprint}:
\begin{itemize}
\item Se comenzó el tratamiento de datos utilizando una librería llamada \textit{Apache Commons CLI}.
\item Se decidió utilizar como motor de plantillas \textit{JinJava}. 
\item En cuanto a las plantillas, y dado que su generación no había sido tan satisfactoria como se pretendía, se decidió junto con el tutor, centrar el trabajo en la plantilla .XML y dejar la plantilla .TEX para un \textit{sprint} futuro.
\end{itemize}

En esta reunión se plantearon nuevos objetivos:
\begin{itemize}
\item Introducción de $JinJava$ para comenzar con la generación de los ficheros deseados.
\item Continuar con la plantilla para ficheros .XML.
\item Establecer una codificación $UTF-8$ para todos los ficheros, tanto en las clases del proyecto como en los ficheros generados.
\end{itemize}

 

\subsection{Sprint 5: 7 abril - 14 abril}
La siguiente reunión tuvo lugar el 7 de abril y se analizaron los progresos conseguidos en el anterior \textit{sprint}. Se desechó el uso de $JinJava$ como motor de plantillas porque no permitía la posibilidad de realizar bucles de repetición. En su lugar, se decidió continuar el trabajo utilizando $Mustache$, ya que cualquier elemento iterable puede ser utilizado. En cuanto a la plantilla para ficheros .XML, no se pudo avanzar lo deseado puesto que aparecieron problemas con $JavaScript$ ya que la plataforma \textit{Moodle} de la Universidad de Burgos no permite su uso.

Los objetivos que se plantearon en esta reunión fueron:
\begin{itemize}
\item Al tener problemas con el uso de $JavaScript$, se plantearon dos posibles soluciones. Una de ellas era utilizar un motor web que, desde Java, permitiese transformar un documento HTML construido de forma dinámica con $JavaScript$, en un documento HTML estático. La otra solución era cambiar todos los script de $JavaScript$ por HTML estático de forma manual. Por tanto, el objetivo era buscar dicho motor web. Si se conseguía encontrar un motor web que transformase los documentos en poco tiempo el objetivo estaría cumplido; sino, se tendría que proceder al cambio de todos los script de forma manual.
\item Buscar información sobre el formato de inicio y fin de las preguntas en Moodle.
\item Durante la reunión se decidió que las gramáticas estuvieran compuestas por letras mayúsculas y letras minúsculas. Las mayúsculas corresponderían a los no terminales y las minúsculas a los terminales. Se estableció que no se utilizasen símbolos como paréntesis, corchetes, llaves, etc. puesto que interfieren con el funcionamiento de la plantilla. De esta manera, se planteó como objetivo cambiar las gramáticas y adaptarlas a los nuevos requisitos planteados en esta reunión.
\end{itemize}



\subsection{Sprint 6: 14 abril - 21 abril}
En la siguiente reunión, que tuvo lugar el 14 de abril, se analizaron los resultados del quinto \textit{sprint}. Debido a la dificultad de encontrar un motor web que permitiese transformar documentos html dinámicos en estáticos, se decidió implantar la segunda solución planteada en la reunión anterior para solucionar el problema que suponía el uso de $JavaScript$. Para ello, se utilizó la propiedad de $Mustache$ que permite utilizar cualquier elemento iterable. Los demás objetivos tuvieron resultados satisfactorios.

En esta reunión se plantearon nuevos objetivos:
\begin{itemize}
\item Crear respuestas alternativas para los símbolos de preanálisis con cierta lógica conociendo la solución correcta.
\item Continuar con la plantilla .XML.
\end{itemize}


\subsection{Sprint 7: 21 abril - 28 abril}
El 21 de abril se realizó una reunión en la que se valoraron los resultados del sexto $sprint$. Todos los resultados fueron satisfactorios.

Ya que todos los objetivos se iban cumpliendo a tiempo, en esta semana se planteó terminar el código para visualizar las tablas del \textit{FIRST} y del \textit{FOLLOW} en una gramática $LL$.


\subsection{Sprint 8: 28 abril - 5 mayo}
En la siguiente reunión celebrada el 28 de abril se comprobó que los objetivos del séptimo \textit{sprint} se habían conseguido, es decir, el código para visualizar las tablas del \textit{FIRST} y del \textit{FOLLOW} en una gramática $LL$ estaba terminado.

En esta reunión se plantearon nuevos objetivos:
\begin{itemize}
\item Se estableció cómo debería ser la plantilla y se programó cambiar la plantilla existente en ese momento para que fuera acorde con lo decidido en la reunión.
\item Crear la tabla de análisis sintáctico predictivo introduciendo la respuesta correcta. Crear un método para obtener las cadenas de múltiples respuestas. Todo ello para la gramática $LL$. 
\end{itemize}



\subsection{Sprint 9: 5 mayo - 12 mayo}
En la reunión del 5 de mayo se valoraron los resultados obtenidos en el octavo $sprint$. Se comprobó que para la gramática $LL$ todo estaba correcto y terminado, menos la tabla de la traza, que se decidió dejar para un futuro $sprint$. Aún así, se decidió meter un parámetro en el prototipo con la cadena para usar posteriormente.

Siguiendo los pasos de la gramática $LL$ se programó para este $sprint$ comenzar a generar las gramáticas $LR$, $SLR$ y $LALR$.
Además, se comentó que se debería comenzar a redactar y escribir las partes de la memoria que ya estuvieran definidas hasta ese momento.



\subsection{Sprint 10: 12 mayo - 19 mayo}
En la siguiente reunión, del 12 de mayo, se comprobó que se habían generado las tablas de conjuntos del análisis $LR$ pero que todavía faltaban por generar las tablas correspondientes a los otros análisis.

En esta reunión se planteó:
\begin{itemize}
\item Crear un método para generar respuestas cortas en Moodle.
\item Crear un método para obtener la tabla de \textit{ACCIÓN} e \textit{IR A}.
\item Seguir trabajando en los análisis restantes.
\item Seguir trabajando en la redacción de la memoria.
\end{itemize}


\subsection{Sprint 11: 19 mayo - 26 mayo}
El 19 de mayo, en la siguiente reunión se comprobó que se habían generado correctamente las respuestas cortas, así como la tabla de \textit{ACCIÓN} e \textit{IR A} en todas los análisis. 
Además se presentó un primer avance de la memoria.
Se comprobó además que el cambio en el posicionamiento de los conjuntos $LR$, $SLR$ y $LALR$ era correcto.

El avance de la aplicación por línea de comando era correcto y quedaban pocas cosas por ir terminando, por tanto, se planteó empezar a desarrollar la interfaz gráfica.
Además, y dado que los análisis estaban prácticamente finalizados, se programó terminarlos creando la tabla para la traza en todas ellas.
Por último, el tutor me facilitó plantillas de memorias de años anteriores para tener modelos a seguir y comenzar con la redacción de los anexos.

\subsection{Sprint 12: 26 mayo - 2 junio}
En la reunión del 26 de mayo se valoraron los resultados del undécimo $sprint$ y se comprobó que todos eran satisfactorios.

Aún avanzando a buen ritmo y con la mayoría del proyecto terminado, se decidió en esta reunión presentarlo en la segunda convocatoria para tener tiempo y terminar de pulir aquellas cosas que todavía no estaban del todo correctas.

En esta reunión se comprobó que, en líneas generales, únicamente quedaba por hacer para la finalización del proyecto y, por tanto, serían los objetivos del decimosegundo \textit{sprint}:
\begin{itemize}
\item Crear pruebas unitarias para evaluar el funcionamiento de los métodos del prototipo.
\item Terminar la interfaz gráfica.
\item Retomar y terminar la plantilla .TEX.
\item Seguir trabajando en la memoria y en sus anexos.
\end{itemize}


\subsection{Sprint 13: 2 junio - 9 junio}
En la siguiente reunión, que tuvo lugar el 2 de junio se comprobaron los avances realizados en el anterior \textit{sprint}. Se comprobó que se habían creado correctamente las pruebas unitarias para evaluar el funcionamiento de los métodos del prototipo. Dichas pruebas se han completado con una cobertura suficiente dadas las características del proyecto. Además, se presentaron los avances conseguidos en la plantilla .TEX para los análisis $LL$ y se plantearon diferentes alternativas para mejorar la forma de la misma.

Como objetivos para el siguiente \textit{sprint}, se plantearon:
\begin{itemize}
\item Continuar con la plantilla .TEX para el resto de análisis e introducir las mejoras planteadas en la reunión.
\item Terminar la interfaz gráfica.
\item Empezar a pensar en la forma y contenido del póster necesario para la defensa del proyecto.
\item Seguir trabajando en la memoria y en sus anexos.
\end{itemize}

\subsection{Sprint 14: 9 junio - 16 junio}

La siguiente reunión tuvo lugar el 9 de junio y se comprobaron los resultados obtenidos en el anterior \textit{sprint}. Así, se presentaron los avances realizados en la interfaz gráfica y un primer borrador de la memoria. Ambos trabajos fueron satisfactorios, aunque se plantearon mejoras en ambos. También se presentaron los avances en la plantilla .TEX para el resto de análisis que, aún sin estar finalizados completamente, avanzaban según los criterios planteados en la anterior reunión.

Por tanto, para el siguiente \textit{sprint} se plantearon los siguientes objetivos:

\begin{itemize}
\item Terminar la plantilla .TEX.
\item Terminar la interfaz gráfica.
\item Corregir los errores de la memoria planteados en la reunión y continuar trabajando en los anexos restantes.
\end{itemize}

\subsection{Sprint 15: 16 junio - 23 junio}

En la siguiente reunión, que tuvo lugar el 16 de junio, se valoraron los avances llevados a cabo en las tareas planteadas para el decimocuarto \textit{sprint}. 

Se presentaron las plantillas .TEX para todos los análisis y se comprobó que eran viables. Aún así, se plantearon mejoras en dichas plantillas que se deberían llevar a cabo en el siguiente \textit{sprint}. 

Se presentaron además los avances llevados a cabo en la interfaz gráfica. En este caso, aunque los avances que se iban realizando eran satisfactorios, el ritmo de avance no era el deseado y se planteó como objetivo para la siguiente semana terminar definitivamente la interfaz.

Finalmente se entregó la memoria y sus anexos corregidos y mejorados con respecto al borrador anterior para su lectura por los tutores.

Por tanto, para el siguiente sprint se planteó como principal objetivo el terminar la interfaz gráfica. También se planteó implementar las mejoras discutidas en las plantillas .TEX.

\subsection{Sprint 15: 23 junio - 30 junio}

La siguiente reunión tuvo lugar el 23 de junio y se valoraron los resultados obtenidos en el \textit{sprint} anterior:

\begin{itemize}
\item Con respecto a la plantilla .TEX, se dieron por buenos los resultados obtenidos.
\item Se plantearon algunas mejoras en la memoria y anexos presentados en la anterior reunión.
\item Con respecto a la interfaz gráfica, se tomaron nuevas decisiones en cuanto al alcance de la misma. Debido a la complejidad de crear una interfaz que funcionara para todos los análisis usando todas las gramáticas, se comprobó que a la hora de exportar los resultados, los datos obtenidos no eran los correctos en algunos de los casos. Así, se tomó la decisión de finalizar la interfaz gráfica llegando solo a los resultados obtenidos hasta ese momento y se planteó como futuros trabajos relacionados con \textit{PLGRAM}, la mejora de la interfaz gráfica.
\end{itemize}

Por tanto, el objetivo planteado para el siguiente \textit{sprint} fue continuar con la documentación del proyecto, planteando tener terminada la memoria y sus anexos para la siguiente reunión.

\subsection{Sprint 16: 30 junio - 7 julio}

En este caso, debido a que en Burgos estaba teniendo lugar la celebración de las Fiestas de San Pedro y San Pablo, no hubo una reunión como tal, sino que se envió la memoria y sus anexos a los tutores para su última revisión antes de la entrega final del proyecto.

\section{Estudio de viabilidad}

El análisis de viabilidad es un requisito imprescindible en cualquier organización de una proyecto, sea éste del tipo que sea.

Son muchos los aspectos que intervienen a la hora de valorar la viabilidad de un proyecto y no es recomendable centrarse únicamente en la relación coste-beneficio del mismo, sino que conviene valorar otros aspectos como la funcionalidad, el tamaño del proyecto, el grado de complejidad del mismo, el equipo de trabajo, etc.

En este caso, y dado la finalidad del proyecto actual, un estudio de viabilidad complejo no es necesario, puesto que el proyecto está planteado dentro de un contexto docente y al amparo de la Universidad de Burgos. 

De todos modos, se plantean algunos aspectos generales, dividiendo el estudio de viabilidad en los dos criterios que se han considerado más importantes: la viabilidad económica y viabilidad legal.

\subsection{Viabilidad económica}

En este apartado se explica el estudio de costes y la viabilidad económica del proyecto, para poder determinar si es viable económicamente o debe desestimarse.

\subsubsection{Estudio de costes}

\begin{enumerate}

\item Costes de personal

Para el cálculo del coste de personal se debe tener en cuenta al desarrollador y a los dos tutores del proyecto que también deben recibir una retribución por el apoyo prestado.

\begin{itemize}
\item Desarrollador: Graduado en Ingeniería Informática que percibirá un salario de 12 $\EUR$/hora. Considerando las horas invertidas a la semana, el desarrollador supone 360$\EUR$/semana. La duración del proyecto ha sido de 16 semanas, por lo que el coste del desarrollador será de 5.760$\EUR$.
\item Tutores: Titulados en Ingeniería Informática que cobrarán 20 $\EUR$/hora. Calculando que la dedicación media va a ser de 1 horas/semana y eniendo en cuenta que la duración del proyecto es de 16 semanas, se tiene:

\textit{2 tutores }$ * \dfrac{\textit{ 20 euros } }{hora}*\dfrac{\textit{1 hora}}{semana}*\textit{16 semanas}=640 \EUR $

\end{itemize}

Sumando el coste del desarrollador y el de los tutores hacen un total de:
5.760$\EUR$ +640$\EUR$ = 6.400$\EUR$.

Al salario bruto hay que añadir el coste de la Seguridad Social (desglosado en contingencias comunes, seguro desempleo y formación profesional) que supone un 30$\%$ aproximadamente sobre
el salario bruto, por lo que dicho coste será: 6.400$\EUR$ * 0,3 = 1.920$\EUR$.

El coste total de personal (sumando el coste de la seguridad social) hace un total de 6.400$\EUR$ + 1.920$\EUR$ = 8.320$\EUR$.

\item Coste de software

Se ha trabajado con herramientas cuya licencia es gratuita, por lo que el coste de software utilizado es de 0$\EUR$.

\item Costes totales

Debido a que los costes de software son 0$\EUR$, el coste total del proyecto asciende a 8.320$\EUR$.
\end{enumerate}

\subsubsection{Análisis de beneficios}

El costal del proyecto ha sido estimado en 8.320$\EUR$ por lo que calculando el precio de la licencia en 150$\EUR$, el número de licencias a vender para obtener beneficios sería de 56 licencias.

Por tanto, para considerar económicamente viable el proyecto, sería necesario vender 56 licencias de \textit{PLGRAM}. Debido al carácter universitario de la aplicación desarrollada, se puede intuir que dicho objetivo es factible de conseguir debido a la cantidad de centros universitarios del país.

\subsection{Viabilidad legal}

En el uso de bibliotecas de terceros, como por ejemplo \textit{Apache Commons Cli}, todo el software producido por la \textit{ASF} (\textit{Apache Software Foundation}) o cualquiera de sus proyectos, está desarrollado bajo los términos de la licencia $Apache$, es decir, la licencia permite al usuario del software la libertad de usar dicho software para cualquier propósito, para distribuirlo, modificarlo y distribuir versiones modificadas del software, bajo los términos de la licencia; sin necesidad de preocuparse por problemas legales futuros.


\apendice{Especificación de Requisitos}

\section{Introducción}
Este anexo tiene como objetivo analizar y documentar las necesidades funcionales mínimas que deben ser soportadas por la aplicación desarrollada. 

Otro objetivo básico del presente anexo es la de priorizar cada una de las funcionalidades a implementar. Si se tiene claro cuáles son las tareas críticas, se podrá llevar a cabo un enfoque más preciso en ellas y asignarlas una mayor cantidad de recursos.

El anexo se estructura presentado una lista de objetivos generales que la aplicación debe cumplir. Son los objetivos que el cliente establece mediante un lenguaje neutral y es trabajo del ingeniero su desglose y formalización. A continuación se incluye un catálogo de los requisitos que deberá cumplir la aplicación.

\section{Objetivos generales}
El desarrollo de la aplicación cumple los siguientes objetivos:
\begin{enumerate}
	\item Construcción de una \textit{GUI} (Interfaz Gráfica de Usuario) que permita visualizar de forma rápida la aplicación desarrollada.
	\item Construcción de una \textit{CLI} (Command Line Interface) que permita trabajar de forma rápida y directa con la aplicación desarrollada.
	\item Procesado y manipulación de gramáticas mediante procesadores del lenguaje.
	\item Obtención del $FIRST$ y del $FOLLOW$ de una gramática.
	\item Obtención de la Tabla de Análisis Sintáctico Predictivo ($TASP$) de una gramática.
	\item Obtención de los conjuntos de items $LR(0)$ y  $LR(1)$.
	\item Obtención de las tablas de análisis sintáctico  \textit{ACCIÓN e IR A} para análisis $SLR(1)$, $LR(1)$ y $LALR(1)$.
	\item Exportación de cuestionarios en formato .XML compatible con la plataforma Moodle.
	
	\item Exportación de cuestionarios en formato \LaTeX{}.
	
	\item Creación de un manual de usuario detallando todas las opciones de la aplicación.
\end{enumerate}

\section{Catálogo de requisitos}
El objetivo de este apartado es definir de forma clara, precisa, completa y verificable todas las funcionalidades y restricciones del sistema a construir.

\subsection{Funciones del producto}
Una de las características más importantes de la aplicación que se ha  desarrollado es su facilidad de uso. 
Dado que su empleo está orientado a convertirse en una herramienta docente, debe ser fácil de usar y posibilitar la creación de cuestionarios de manera rápida, sencilla y comprensible.

Como se explica en el \textit{Anexo E: Manual de usuario}, los archivos obtenidos de la aplicación tendrán extensiones .TEX y .XML, tratando de cubrir un mayor abanico de posibilidades de uso de la aplicación y dando una mayor cobertura a la hora de obtener los cuestionarios.

En el caso de los archivos .XML es necesario el uso de la plataforma \textit{Moodle}, que al ser internacional, se puede utilizar en varios idiomas, dándole así una mayor difusión a la aplicación.

\subsection{Requisitos de usuario}
El usuario debe ser capaz de utilizar la aplicación mediante la Interfaz Gráfica o mediante la Consola del sistema. 

La ayuda será un punto a tener en cuenta en el caso de utilizar la aplicación mediante la consola del sistema.
 
En todo momento el usuario podrá consultar el \textit{Manual del Usuario} para
orientarse y resolver dudas que le puedan surgir durante el manejo de la misma.

\subsection{Requisitos del sistema}
Los requisitos mínimos del hardware del sistema serán los mismos que para ejecutar java:


\url{https://www.java.com/es/download/help/sysreq.xml}

Los requisitos software necesarios son:
\begin{itemize}
\item Sistema operativo \textit{Windows}, \textit{Linux} o \textit{Macintosh}.
\item Máquina virtual de java ($JDK 8$).
\item Consola de sistema operativo capaz de interpretar Unicode, solo para la interfaz en línea de comandos. Si no soportara Unicode, las tildes u otros símbolos podrían ser mal representados.
\end{itemize}

\section{Especificación de requisitos}
En este apartado se realizará un análisis de los requisitos de diseño de la aplicación. Para ello, se detallarán los requisitos funcionales y no funcionales de la aplicación.

\subsection{Requisitos funcionales}
La aplicación debe ser capaz de cumplir los siguientes requisitos:

\begin{itemize}
\item \textit{RF1: Cargar gramática y análisis}. Permite al usuario seleccionar una gramática disponible en la carpeta \textit{gramaticas} y un tipo de análisis.
\item \textit{RF2: Obtener análisis}. Se obtienen los datos del análisis solicitado para la gramática seleccionada.
\item \textit{RF3: Crear cadenas}. Se obtienen las cadenas de los datos del análisis.
\item \textit{RF4: Obtener fichero} .TEX. Se utiliza la plantilla para obtener el fichero .TEX.
\item \textit{RF5: Obtener fichero} .XML. Se utiliza la plantilla para obtener el fichero .XML.
\end{itemize}

\imagen{diagramaprincipal}{Diagrama principal de casos de uso.}


\subsection{Requisitos no funcionales}
A continuación se describen todos aquellos requisitos no funcionales importantes para considerar en el diseño:


\begin{itemize}
\item \textit{RNF1: Eficiencia}. Se debe buscar la eficiencia de la aplicación para que pueda ser usada con gramáticas complejas.
\item \textit{RNF2: Extensibilidad}. Se pueden definir nuevas capacidades en la aplicación mediante su extensión a otro tipo de ficheros.
\end{itemize}
\apendice{Especificación de diseño}

\section{Introducción}

\section{Diseño de datos}

\section{Diseño procedimental}

\section{Diseño arquitectónico}



\apendice{Documentación técnica de programación}

\section{Introducción}
En este capítulo se detallan los aspectos más relevantes de la implementación del diseño expuesto en el apéndice C. Además se detallarán las convenciones de estilo y codificación utilizadas durante la fase de programación.

El desarrollo se ha realizado en la versión 8 de \textit{Java}. Para la implementación del núcleo de la aplicación se ha utilizado \textit{Eclipse} como entorno integrado de desarrollo. Para el desarrollo de
las interfaces de usuario se ha empleado \textit{Swing} y para el control de versiones se ha utilizado \textit{VersionOne}.

El motivo de escoger \textit{Eclipse} fue que ya se había utilizado en proyectos anteriores, conociendo así sus puntos fuertes y débiles. La potencia de esta herramienta facilita las tareas de programación, sobre todo a la hora de refactorizar o depurar el código fuente. Otra de las ventajas que ofrece, es la cantidad de plugins disponibles: control de tareas, control de versiones, pruebas, etc.

Como en todo proyecto software, era necesario disponer de un control de versiones potente y versátil para controlar los fuentes desarrollados(\textit{Github}) y un control de tareas. Para esta tarea se barajaron diversos programas: \textit{Subversion}, \textit{GIT} y \textit{Plastic SCM}, entre otros. Se escogió \textit{VersionOne}. La licencia gratuita de \textit{VersionOne} caducaba a los 30 días de comenzar a usarlo, borrando todos los datos almacenados en ese momento. Por ello, cuando el proyecto estaba ya muy avanzado, se perdió toda la información que se había ido almacenando en \textit{VersionOne} sin que hubiera ninguna posibilidad de recuperarla\footnote{\url {https://www53.v1host.com/Ubu44/Default.aspx?menu=MyHomeEnterpriseGettingStartedPage}}.

Para la realización de las pruebas unitarias se ha optado por \textit{JUnit}. Para las pruebas de cobertura se ha utilizado un plugin de \textit{Eclipse} llamado \textit{EclEmma} por la calidad en el desempeño de sus tareas, además de ser un plugin ya conocido al haberlo utilizado durante la carrera.

Para la obtención de métricas se ha utilizado \textit{SourceMonitor} por su facilidad de
uso y porque, entre otras cosas, aportaba las métricas necesarias para los informes.

\section{Estructura de directorios}

La estructura de directorios que se ha seguido en el soporte físico distribuido junto a la documentación, es la siguiente:
\begin{itemize}
\item PLGRAM: Este directorio contiene la aplicación desarrollada, incluyendo los fuentes, las plantillas, los ejecutables y las pruebas. Además, incluye el \textit{build.xml} para \textit{Ant}.
\item Memoria: Este directorio contiene la memoria del trabajo.
\item Documentación técnica: Este directorio contiene la documentación que se está leyendo en este momento.
\end{itemize}

\section{Manual del programador}

El manual del programador se encuentra dividido en dos secciones: en la primera, aparecen detalladas las herramientas utilizadas durante el desarrollo y, en la segunda, el fichero de configuración
de Ant explicando cada uno de los objetivos (o \textit{target}).

\subsection{Herramientas utilizadas}

Las herramientas utilizadas han sido cuidadosamente explicadas en la memoria, en el capitulo 4 Técnicas y herramientas, sección 2 Herramientas.


\subsection{Configuración de la construcción}



Para la construcción de los ejecutables, pruebas, métricas y demás objetos se ha utilizado \textit{Ant}. Esta herramienta sería el equivalente en \textit{Java} al \textit{Make} de \textit{C}.

\textit{Ant} utiliza ficheros de configuración escritos en XML. Para el proyecto, ha sido generado el fichero \textit{build.xml}, que se encuentra en el directorio PLGRAM.

\section{Descarga, instalación y ejecución del proyecto}
Todo el código de la aplicación está detalladamente comentado para facilitar a un futuro programador la ampliación y programación que se desee realizar.

\subsection{Descarga}
El proyecto puede obtenerse directamente o a través de su repositorio de \textit{Github}\footnote{\url {https://github.com/vrt0004/Trabajo}}, que ofrece dos opciones:

\begin{itemize}
\item Puede descargarse como un fichero comprimido mediante la opción Download  ZIP.
\item Puede descargarse utilizando  uno  de  los  clientes  propietarios  de  Github  con  la opción Clone  in  Desktop.
\end{itemize}

\subsection{Instalación y ejecución}
La aplicación se distribuye de dos formas diferentes: $PLGRAM$ y $PLGRAMlineCommand$.

\begin{itemize}
\item Para ejecutar $PLGRAM$ se dispone de un fichero jar ejecutable.
\imagen{iconoPLGRAM}{Icono PLGRAM.}
\item Para ejecutar $PLGRAMlineCommand$ se debe acceder mediante la línea de comandos a la carpeta que contiene el fichero.

Esta carpeta debe contener también una carpeta llamada \textit{gramaticas} donde estarán las gramáticas guardadas. Además, también estarán en ese directorio las plantillas necesarias \textit{(plantillaTEX y plantillamoodle)}.
\end{itemize}
\newpage
\section{Pruebas del sistema}

La aplicación utiliza JUnit 4 para realizar sus pruebas, y cuenta con un total de 27 pruebas. Las pruebas excluyen las clases autogeneradas (como por ejemplo las pertenecientes al paquete parser, generadas porJavaCC). Las herramientas aplicadas indican una cobertura general superior al $55\%$, con más del $90\%$ en las clases más relevantes.
\vspace{0.5cm}
\imagen{cobertura}{Resultados del análisis de cobertura de pruebas.}
\newpage
\section{Métricas}

Para las distintas métricas se ha utilizado \textit{SourceMonitor}.

El resultado del análisis obtenido con \textit{SourceMonitor} se puede observar en la tabla siguiente. Se muestra también un gráfico de Kiviat obtenido con \textit{SourceMonitor}.

En ambos, aparecen los umbrales definidos por la Universidad de Burgos. Los valores que no se encuentran dentro del umbral establecido han sido resaltados en negrita.

\imagen {metricas}{Métricas obtenidas con SourceMonitor.}

\begin{tabular}{|l|c|c|c|}
\hline Métrica & Resultado &  \begin{tabular}[c]{@{}c@{}}Mín.\\ UBU\end{tabular} & \begin{tabular}[c]{@{}c@{}}Máx.\\ UBU\end{tabular}\\





\hline Líneas de código & 10.506 & &\\
\hline Número de sentencias & 5.966 &&\\
\hline Porcentaje de sentencias condicionales& 20,1 &&\\
\hline Número de llamadas a métodos &4.124 &&\\
\hline \textbf{Porcentaje de líneas de comentarios} & 25,3& 8 &20\\
\hline Número de clases e interfaces &77 &&\\
\hline Número de métodos por clase &6,03& 4 &16\\
\hline Media de sentencias por método& 10,21& 6& 12\\
\hline \textbf{Complejidad máxima} & 9+ &2& 8\\
\hline Media de complejidad &2,78 & 2& 4\\
\hline \textbf{Máxima profundidad en bloques} &9+ & 3 &7\\
\hline \textbf{Media de profundidad de bloques} & 2,96 & 1 & 2,2\\
 \hline
\end{tabular}


Aunque la complejidad máxima y la profundidad máxima se salen de los valores establecidos, la media de ambas métricas se encuentra en el umbral o muy cerca de estarlo.

\newpage
A continuación se muestra un diagrama de kiviatt de la aplicación PLGRAM.

Los diagramas de Kiviatt son unos gráficos circulares en cuyos ejes radiales se representan diferentes indices de prestaciones. Las intersecciones entre los radios y la circunferencia representan los valores máximos que pueden alcanzar las variables representadas en los mismos. Aunque en principio, el número de ejes que puede tener un gráfico de este tipo es arbitrario y depende de los datos que se van a representar, se suelen seguir unos convenios de representación.
\imagen{kiviat}{Gráfico de Kiviat obtenido con SourceMonitor.}
\apendice{Documentación de usuario}

\section{Introducción}
En esta sección se explica cómo llevar a cabo la instalación de las herramientas implementadas en el presente proyecto y los requisitos del sistema.
\section{Requisitos de usuarios}
Los requisitos mínimos del hardware del sistema serán los mismos que para ejecutar java\cite{requisitos}

Los requisitos software necesarios son:

\begin{itemize}
\item Sistema operativo \textit{Windows}, \textit{Linux} o \textit{Macintosh}.
\item Máquina virtual de java ($JDK 8$).
\item Consola de sistema operativo capaz de interpretar Unicode, solo para la interfaz en línea de comandos. Si no soportara Unicode, las tildes u otros símbolos podrían ser mal representados.
\end{itemize}
\section{Instalación}
La aplicación se distribuye de dos formas diferentes: $PLGRAM$ y $PLGRAMlineCommand$:

\begin{itemize}
\item Para ejecutar $PLGRAM$ se dispone de un fichero jar ejecutable.
\imagen{iconoPLGRAM}{Icono PLGRAM.}
\item Para ejecutar $PLGRAMlineCommand$ se debe acceder mediante la línea de comandos a la carpeta que contiene el fichero.

Esta carpeta debe contener también una carpeta llamada \textit{gramaticas} donde estarán las gramáticas guardadas. Además, también estarán en ese directorio las plantillas necesarias \textit{(plantillaTEX y plantillamoodle)}.


\end{itemize}


\section{Manual del usuario}

\subsection{Aplicación línea de comandos}

Para ejecutar la aplicación se debe escribir el siguiente comando:
\begin{verbatim}
 java -jar PLGRAMlineCommand.jar -g -t [-i] [-ca] [-o] 
\end{verbatim}

 
 Los diferentes argumentos son:
 \begin{itemize}
 	\item $-g$ Nombre de la gramática a analizar (que debe estar en la carpeta \textit{gramaticas} mencionada anteriormente).
 	\item $-t$ Tipo de análisis a realizar $[LL,LR,SLR,LALR].$
 	\item $-i$ Extensión del informe $[TEX, XML, ALL]$.
 	\item $-ca$ Cadena a analizar.
 	\item $-o$ Nombre del fichero de salida.
 	
\end{itemize} 



La gramática $-g$ y el tipo de análisis a realizar $-t$ son obligatorios. La gramática debe estar en la carpeta \textit{gramaticas}.
 
El parámetro $-t$ tiene diferentes opciones:
  \begin{itemize}
 	 \item $LL$ Para realizar un análisis $LL$.
 	 \item $LR$ Para realizar un análisis $LR$.
 	 \item $SLR$ Para realizar un análisis $SLR$.
 	 \item $LALR$ Para realizar un análisis $LALR$.
 \end{itemize}
  
 
Si no se introduce alguno de los parámetros obligatorios (\textit{-g} , \textit{-t}) o el parámetro es incorrecto se indicará que falta el parámetro.
 
Así, si se introduce la siguiente línea de código 
\begin{verbatim}
 java -jar PLGRAMlineCommand.jar  
\end{verbatim}

La aplicación requerirá la gramática, tal y como se muestra a continuación:

\imagen{consola2}{$PLGRAM$ sin argumento gramática.}

La primera vez que se ejecuta la aplicación y se genera un fichero se crea, en el directorio donde se encuentra el ejecutable, una nueva carpeta llamada \textit{informes}, donde se irán guardando todos los ficheros generados posteriormente.

Si solo se introducen los parámetros de la gramática y del tipo de análisis que se quiere realizar, se ejecuta el programa con el resto de valores por defecto:
\begin{itemize}
\item $-i$ con valor $ALL$.
\item $-o$ que se denominará de manera genérica como el tipo de análisis elegido concatenado al nombre de la gramática seleccionada. Por ejemplo, para un análisis $LL$ y una gramática $G3$, los ficheros generados se denominarán $LLG3$, cada uno con su extensión correspondiente.
\item $-ca$ no se generará ninguna tabla de primeros pasos de la traza de análisis.
\end{itemize}


En el supuesto de que se quiera modificar alguno de los valores que por defecto se les asignan a los parámetros anteriores, bastará con introducir en cada uno de ellos los argumentos deseados:
\begin{itemize}
\item $-i$ se podrá elegir la extensión del fichero de salida, que será XML en el caso de querer obtener un cuestionario virtual o TEX en el caso de que se quiera para una futura impresión.
\item $-o$ se podrá nombrar los ficheros de salida según se desee.
\item $-ca$ se podrá introducir una cadena para  que la aplicación genere una tabla de primeros pasos de la traza de análisis de dicha cadena.
La cadena deberá introducirse entrecomillado y cada item separado por un espacio. Por ejemplo, $-ca$ ``$item1$ $item2$ $item3$ $item4$".
\end{itemize}


Una vez introducidos todos los parámetros, se muestra la información de la gramática elegida a través de la consola.

A continuación se muestran los datos obtenidos al introducir una gramática $gramatica1.yc$ y un tipo de análisis $LL$. El resto de parámetros se han dejado por defecto.

\imagen{consola4}{$PLGRAM$ con argumentos gramática y tipo de análisis.}

Como se puede observar, al final de la información obtenida, la consola nos muestra que los ficheros se han generado correctamente y su ruta.


\subsection{Aplicación interfaz gráfica}
La interfaz gráfica pretende que el usuario tenga una visión completa y sencilla de la aplicación.
\imagen{gui1}{Aplicación \textit{PLGRAM}.}
\newpage
\subsubsection{Añadir preguntas}
Para añadir preguntas mediante la interfaz gráfica, se tienen dos opciones:

\begin{itemize}
\item Desde el menú Archivo. En él se ofrecen la opción de añadir preguntas utilizando los cuatro tipos de análisis.
\imagen{gui2}{Añadir preguntas desde menú Archivo.}
\item Utilizando el desplegable de la parte inferior de la aplicación. Al abrir dicho desplegable, se nos muestran los cuatro tipos de análisis que se pueden elegir a la hora de añadir una pregunta. Una vez seleccionado el tipo de análisis, bastaría con pulsar el botón identificado con el símbolo '+' situado a la izquierda del desplegable.
\imagen{gui3}{Añadir preguntas desde el desplegable.}
\end{itemize}

Utilizando cualquiera de las dos opciones anteriores, el resultado obtenido es una panel de preguntas que aparecerá en la parte superior izquierda de la aplicación.
\imagen{gui4}{Panel de preguntas.}

\subsubsection{Borrar preguntas}
En el caso de querer borrar una pregunta se realiza mediante el botón situado a la izquierda del panel de preguntas identificado con el símbolo '-'.
\imagen{gui5}{Borrar preguntas.}

\subsubsection{Cambiar el orden de las preguntas}
Si se tienen varias preguntas abiertas a la vez, los controles situados a la izquierda del panel de preguntas permiten modificar la posición de las preguntas hacia arriba y hacia abajo dentro de la ventana izquieda.
\imagen{gui6}{Cambiar el orden de las preguntas.}

\subsubsection{Cargar gramáticas}
Una vez elegido el tipo de análisis, se permite, desde el panel de preguntas, cargar la gramática deseada. Para ello, habría que pulsar el botón Cargar de dicho panel.
\imagen{gui7}{Cargar gramáticas.}

Aparece un explorador de ficheros en el que se puede seleccionar una de las gramáticas que se tengan guardadas.

Una vez abierta la gramática, el panel de preguntas ser rellenará con los datos de dicha gramática.
\imagen{gui8}{Ejemplo de una gramática cargada en el panel de preguntas.}
\newpage
\subsubsection{Añadir cadena}
La aplicación permite añadir una cadena si se desea generar posteriormente su traza.Para ello se deberá introducir en el cuadro habilitado una cadena con los elementos separados por espacios. 

\imagen{gui10}{Ejemplo de cadena.}
\newpage
\subsubsection{Mostrar gramáticas}
Elegido el tipo de análisis, la gramática deseada e introducida o no una cadena, pulsando en el botón Mostrar situado en la parte inferior del panel de preguntas, se previsualiza el resultado del análisis en la parte derecha de la pantalla.
\imagen{gui11}{Ejemplo de la previsualización de un análisis.}

Cada panel de preguntas generado para cada una de las preguntas deseadas tiene su propio botón mostrar. 
\newpage
\subsubsection{Exportar ficheros}
Una vez obtenida la previsualización del análisis de la gramática, la aplicación permite exportar los resultados a dos tipos de ficheros: ficheros .XML y ficheros .TEX.

Para ello, debe elegirse la opción que se desee dentro del menú Exportar.
\imagen{gui12}{Exportar ficheros.}

En cualquiera de las dos opciones, aparece un cuadro de diálogo que nos informa de que el fichero se ha generado correctamente. Dicho fichero se guarda en el directorio donde se esté ejecutando la aplicación \textit{PLGRAM.jar}.
\newpage
\subsubsection{Página web}
En el menú Ayuda se encuentra la opción Página web que nos redirecciona a la página web del proyecto.

\subsubsection{Acerca de}
En el menú Ayuda se encuentra la información de la aplicación.
\imagen{gui13}{Acerca de.}

\subsection{Ficheros obtenidos}

Bien a partir de la \textit{GUI} o mediente la \textit{CLI} se han podido generar ficheros en los formatos .TEX o .XML
\newpage
\subsubsection{Fichero .XML}

El fichero .XML obtenido tiene una estructura general común sea cual sea la gramática y el tipo de análisis realizado. Esta estructura sigue el siguiente esquema general:
\imagen{EstructuraXML}{Estructura del fichero .XML.}

Para poder utilizar el fichero .XML obtenido es necesario poder acceder a la plataforma de aprendizaje de Moodle.

Los pasos que se deben seguir son:
\begin{enumerate}
\item Ir a la página principal del curso.
\item Hacer clic en Banco de preguntas desde el bloque de Administración.
\item Pulsar en Importar archivos.
\item Seleccionar la opción de formato XML Moodle.
\item Desde la siguiente ventana pulsar en Examinar.
\item Seleccionar el fichero .XML en nuestro ordenador.
\item Pulsar en Importar este archivo.
\end{enumerate}

Los puntos 5 y 6 se pueden realizar arrastrando el archivo al recuadro que se muestra.


A la hora de importar el fichero .XML hay que tener en cuenta lo siguiente:
\begin{itemize}
\item Los nombres de los archivos y carpetas no deben tener caracteres especiales como palabras acentuadas, tabulaciones, ñ, retornos de carro, espacios en blanco, ni símbolos de sistema como / o : , además no es recomendable combinar mayúsculas y minúsculas en los nombres.
\item El tamaño del archivo no debe superar el límite permitido que se muestra en la ventana.

\imagen{Moodle1}{Importar fichero .XML en Moodle.}
\end{itemize} 

Una vez importado el fichero .XML deseado, aparece una primera previsualización de los datos que se obtendrán a continuación. Basta con dar a continuar para acceder a la siguiente pantalla.

En ella, aparecen todos los ficheros que se hayan importado. Para visualizar uno de ellos, se hace click en la lupa del archivo correspondiente.
\imagen{moodle2}{Lista de ficheros .XML en Moodle.}

Un ejemplo de cuestionario obtenido sería el siguiente:
\imagen{moodle3}{Ejemplo de cuestionario \textit{LL} resuelto en Moodle.}
\imagen{moodle4}{Ejemplo de cuestionario  \textit{LR}(1)  resuelto en Moodle.Conjunto de items}
\imagen{moodle5}{Ejemplo de cuestionario  \textit{LR}(1) resuelto en Moodle.Tabla de ACCIÓN e IR A}
\subsubsection{Fichero .TEX}

Al tratarse de un fichero .TEX se requiere un editor de textos adaptado a \LaTeX{} para utilizarlo.

Una vez abierto el fichero, su tratamiento es muy sencillo. Bastará con compilar el código y se obtendrá directamente un visionado del cuestionario que se ha realizado.

Se muestra un ejemplo de un fichero .TEX abierto con \TeX{}MaKer. De manera general, se puede observar el código a la izquierda y del documento generado a la derecha.

\imagen{ejemplotex1}{Ejemplo fichero .TEX para un análisis \textit{LL}.}

Para una mejor comprensión, se muestra un detalle de cómo compilar el código para obtener el documento. Para ello, basta con hacer click en las flechas resaltadas en la parte superior de la pantalla del editor.

\imagen{ejemplotex2}{Compilación de un fichero .TEX.}

Es necesario destacar que para cada tipo de análisis elegido, se genera un fichero .TEX diferente siguiendo los requisitos de cada una de ellas. 

La aplicación ofrece dos opciones en lo que a los ficheros .TEX respecta. Así, se podrá obtener un documento en el que el cuestionario esté completo con las respuestas correctas y otro documento en el que el cuestionario se encuentre vacío, pensado para ser respondido por un tercero.

\begin{itemize}
\item Obtención del cuestionario completo. Es la opción que se genera por defecto y por tanto no será necesario tratar el código para su obtención. De cualquier manera, debería estar activa la siguiente línea:
\begin{verbatim}
\newcommand{\h}[1]{#1}
\end{verbatim}
\imagen{ejemplotex3}{Obtención del cuestionario completo.}
\item Obtención del cuestionario en blanco. Para ello, bastaría con eliminar el símbolo $\%$ de delante de la siguiente línea:

\begin{verbatim}
%\renewcommand{\h}[1]
\end{verbatim}

\imagen{ejemplotex4}{Obtención del cuestionario en blanco: descomentar.}
Así, ambas líneas deberían estar activas.
\imagen{ejemplotex5}{Obtención del cuestionario en blanco.}

\end{itemize}

Siempre que se compile código, se crea un archivo .PDF en el mismo directorio donde se encuentra el archivo .TEX, con los cuestionarios elegidos.



\bibliographystyle{plain}
\bibliography{bibliografiaAnexos}

\end{document}